\textbf{Enunciado do Problema:} \\

Seja $T$ o operador linear de $\mathbb{R}^4$ cuja matriz em relação à base canônica é

\[
A = \begin{pmatrix}
1 & 3 & 3 & 0 \\
3 & 1 & 3 & 0 \\
-3 & -3 & -5 & 0 \\
4 & -2 & 3 & 1
\end{pmatrix}
\]

Determine operadores lineares $D$ e $N$ de $\mathbb{R}^4$ tais que $T = D + N$, $D$ é diagonalizável, $N$ é nilpotente e $DN = ND$.

\bigskip

\textbf{Resolução Detalhada} \\

A construção da decomposição $T = D + N$ é realizada através da Forma Canônica de Jordan do operador $T$. Se $J$ é a forma de Jordan de $A$, com $A = PJP^{-1}$, então a decomposição de Jordan de $J$ em sua parte diagonal $J_D$ e sua parte nilpotente $J_N$ (estritamente triangular superior), $J = J_D + J_N$, nos dá a decomposição de $A$:

\[
A = P(J_D + J_N)P^{-1} = \underbrace{PJ_DP^{-1}}_{D} + \underbrace{PJ_NP^{-1}}_{N}
\]

A matriz $D$ será a representação de $D$, e $N$ a representação de $N$.

\bigskip

\textbf{Passo 1: Polinômio Característico e Autovalores} \\

Primeiro, encontramos os autovalores de $T$ calculando o polinômio característico $p_T(\lambda) = \det(A - \lambda I)$.

\[
p_T(\lambda) = \det \begin{pmatrix}
1-\lambda & 3 & 3 & 0 \\
3 & 1-\lambda & 3 & 0 \\
-3 & -3 & -5-\lambda & 0 \\
4 & -2 & 3 & 1-\lambda
\end{pmatrix}
\]

Expandindo ao longo da quarta coluna:

\[
p_T(\lambda) = (1-\lambda)\det \begin{pmatrix}
1-\lambda & 3 & 3 \\
3 & 1-\lambda & 3 \\
-3 & -3 & -5-\lambda
\end{pmatrix}
\]

O determinante $3 \times 3$ é:
\[
(1-\lambda)((1-\lambda)(-5-\lambda)+9) - 3(3(-5-\lambda)+9) + 3(-9+3(1-\lambda))
\]
\[
= (1-\lambda)(\lambda^2+4\lambda+4) - 3(-6-3\lambda) + 3(-6-3\lambda)
\]
\[
= (1-\lambda)(\lambda+2)^2
\]

Portanto,
\[
p_T(\lambda) = (1-\lambda)(\lambda+2)^2 = (1-\lambda)(\lambda+2)^2
\]

Os autovalores são $\lambda_1 = 1$ (multiplicidade algébrica 2) e $\lambda_2 = -2$ (multiplicidade algébrica 2).

\bigskip

\textbf{Passo 2: Autoespaços e Diagonalizabilidade} \\

\textbf{Para $\lambda_1 = 1$:} O autoespaço $E_1 = \ker(A - I)$. A matriz $A - I$ é:

\[
\begin{pmatrix}
0 & 3 & 3 & 0 \\
3 & 0 & 3 & 0 \\
-3 & -3 & -6 & 0 \\
4 & -2 & 3 & 0
\end{pmatrix}
\]

\[
\longrightarrow \text{escalonamento} \longrightarrow
\begin{pmatrix}
1 & 0 & 0 & 0 \\
0 & 1 & 0 & 0 \\
0 & 0 & 1 & 0 \\
0 & 0 & 0 & 0
\end{pmatrix}
\]

A solução para $(A-I)v=0$ é $x=y=z=0$. O autoespaço é gerado por $v_1 = (0,0,0,1)^T$.

A multiplicidade geométrica de $\lambda_1 = 1$ é $\dim(E_1) = 1$, que é menor que a multiplicidade algébrica (2). Portanto, \textbf{o operador T não é diagonalizável.}

\bigskip

\textbf{Para $\lambda_2 = -2$:} O autoespaço $E_{-2} = \ker(A+2I)$. A matriz $A+2I$ é:

\[
\begin{pmatrix}
3 & 3 & 3 & 0 \\
3 & 3 & 3 & 0 \\
-3 & -3 & -3 & 0 \\
4 & -2 & 3 & 3
\end{pmatrix}
\]

As três primeiras linhas nos dão $x+y+z=0$. A última nos dá $4x-2y+3z+3w=0$. Resolvendo o sistema, encontramos soluções da forma $y(5,1,-6,0)^T + w(-3,0,3,1)^T$.

O autoespaço $E_{-2}$ é gerado pelos autovetores $v_2=(5,1,-6,0)^T$ e $v_3=(-3,0,3,1)^T$. A multiplicidade geométrica de $\lambda_2=-2$ é $\dim(E_{-2})=2$, que é igual à sua multiplicidade algébrica.

\bigskip

\textbf{Passo 3: Base de Jordan e Matriz de Passagem $P$} \\

A Forma de Jordan $J$ terá um bloco de Jordan $2\times2$ para $\lambda_1=1$ e dois blocos $1\times1$ para $\lambda_2=-2$.

\[
J = \begin{pmatrix}
1 & 1 & 0 & 0 \\
0 & 1 & 0 & 0 \\
0 & 0 & -2 & 0 \\
0 & 0 & 0 & -2
\end{pmatrix}
\]

Para construir a matriz de passagem $P$, precisamos de uma base de Jordan $\{u_1,u_2,v_2,v_3\}$.

- $v_2, v_3$ são os autovetores para $\lambda_2=-2$ que já encontramos.
- Para $\lambda_1=1$, precisamos de um autovetor generalizado $u_1$ e um autovetor $u_2$ tais que $(A-I)u_1=u_2$ e $(A-I)u_2=0$.

Primeiro, encontramos o autoespaço generalizado $K_1 = \ker((A-I)^2)$.

\[
(A-I)^2 = \begin{pmatrix}
0 & -9 & -9 & 0 \\
-9 & 0 & -9 & 0 \\
9 & 9 & 18 & 0 \\
-15 & 3 & -12 & 0
\end{pmatrix}
\]

Resolvendo $(A-I)^2v=0$, obtemos as condições $x=-z$ e $y=-z$. Uma base para $K_1$ é $\{(-1,-1,1,0)^T,(0,0,0,1)^T\}$.

Escolhemos um vetor $u_1\in K_1$ que não esteja em $E_1$. Seja $u_1=(-1,-1,1,0)^T$. Calculamos $u_2=(A-I)u_1$:

\[
u_2 = \begin{pmatrix}
0 & 3 & 3 & 0 \\
3 & 0 & 3 & 0 \\
-3 & -3 & -6 & 0 \\
4 & -2 & 3 & 0
\end{pmatrix}
\begin{pmatrix}
-1 \\ -1 \\ 1 \\ 0
\end{pmatrix}
= \begin{pmatrix}
0 \\ 0 \\ 0 \\ 1
\end{pmatrix}
\]

Este $u_2$ está em $E_1$, como esperado. Nossa base de Jordan é $\{u_1,u_2,v_2,v_3\}$. A matriz de passagem $P$ tem estes vetores como colunas:

\[
P = \begin{pmatrix}
-1 & 0 & 5 & -3 \\
-1 & 0 & 1 & 0 \\
1 & 0 & -6 & 3 \\
0 & 1 & 0 & 1
\end{pmatrix}
\]

\bigskip

\textbf{Passo 4: Cálculo de $P^{-1}$ e das Matrizes $D$ e $N$} \\

O cálculo da inversa de $P$ é um exercício de eliminação de Gauss-Jordan sobre $[P|I]$. O resultado é:

\[
P^{-1} = \tfrac{1}{4}
\begin{pmatrix}
3 & -5 & -3 & 1 \\
1 & -5 & -3 & 1 \\
1 & -1 & 1 & -1 \\
-1 & 5 & 3 & -1
\end{pmatrix}
\]

As partes diagonal e nilpotente de $J$ são:

\[
J_D = \begin{pmatrix}
1 & 0 & 0 & 0 \\
0 & 1 & 0 & 0 \\
0 & 0 & -2 & 0 \\
0 & 0 & 0 & -2
\end{pmatrix}, \quad
J_N = \begin{pmatrix}
0 & 1 & 0 & 0 \\
0 & 0 & 0 & 0 \\
0 & 0 & 0 & 0 \\
0 & 0 & 0 & 0
\end{pmatrix}
\]

Calculamos $N = PJ_NP^{-1}$:

\[
(PJ_N)P^{-1} = \tfrac{1}{4}
\begin{pmatrix}
0 & -1 & 0 & 0 \\
0 & -1 & 0 & 0 \\
0 & 1 & 0 & 0 \\
0 & 0 & 0 & 0
\end{pmatrix}
\begin{pmatrix}
3 & -5 & -3 & 1 \\
1 & -5 & -3 & 1 \\
1 & -1 & 1 & -1 \\
-1 & 5 & 3 & -1
\end{pmatrix}
= \tfrac{1}{4}
\begin{pmatrix}
-1 & 5 & 3 & -1 \\
-1 & 5 & 3 & -1 \\
1 & -5 & -3 & 1 \\
0 & 0 & 0 & 0
\end{pmatrix}
\]

Esta é a matriz do operador nilpotente $N$. Pode-se verificar que $N^2=0$.

Finalmente, encontramos $D$ através de $D=A-N$:

\[
D = \begin{pmatrix}
1 & 3 & 3 & 0 \\
3 & 1 & 3 & 0 \\
-3 & -3 & -5 & 0 \\
4 & -2 & 3 & 1
\end{pmatrix}
- \tfrac{1}{4}
\begin{pmatrix}
-1 & 5 & 3 & -1 \\
-1 & 5 & 3 & -1 \\
1 & -5 & -3 & 1 \\
0 & 0 & 0 & 0
\end{pmatrix}
\]

\[
= \tfrac{1}{4}\begin{pmatrix}
5 & 7 & 9 & 1 \\
13 & -1 & 9 & 1 \\
-13 & -7 & -17 & -1 \\
16 & -8 & 12 & 4
\end{pmatrix}
\]

\bigskip

\textbf{Conclusão} \\

A decomposição de Jordan-Chevalley do operador $T$ é $T = D+N$, onde os operadores $D$ e $N$ são representados na base canônica pelas matrizes:

\[
D = \tfrac{1}{4}\begin{pmatrix}
5 & 7 & 9 & 1 \\
13 & -1 & 9 & 1 \\
-13 & -7 & -17 & -1 \\
16 & -8 & 12 & 4
\end{pmatrix},
\quad
N = \tfrac{1}{4}\begin{pmatrix}
-1 & 5 & 3 & -1 \\
-1 & 5 & 3 & -1 \\
1 & -5 & -3 & 1 \\
0 & 0 & 0 & 0
\end{pmatrix}
\]

A construção através da Forma de Jordan garante, por teorema, que $D$ é diagonalizável (com autovalores 1 e -2), $N$ é nilpotente (especificamente, $N^2=0$), e que eles comutam ($DN=ND$).